% !TeX root = manuscrito.tex
% Plantilla vacía del nuevo artículo. Compilar con ../compilar-articulo.sh.
\documentclass[pdflatex,sn-mathphys-num]{sn-jnl}

\usepackage[utf8]{inputenc}
\usepackage[T1]{fontenc}
\usepackage[spanish,es-noshorthands,es-tabla]{babel}
\usepackage{graphicx}
\usepackage{multirow}
\usepackage{amsmath,amssymb,amsfonts}
\usepackage{amsthm}
\usepackage{mathrsfs}
\usepackage[title]{appendix}
\usepackage{xcolor}
\usepackage{textcomp}
\usepackage{manyfoot}
\usepackage{booktabs}
\usepackage{array,longtable,tabularx}
\usepackage{tikz}
\graphicspath{{figuras/}}

\begin{document}

\title{Título del nuevo artículo}
\author{\fnm{Nombre} \sur{Apellido}}
\maketitle

\section{Introducción}
% Redactar aquí el nuevo artículo.

% Activar cuando se agreguen las primeras citas del nuevo artículo:
% \bibliography{referencias}

\end{document}
